% 第 7 章 IC Validator Explorer LVS
\bichapter{IC Validator Explorer LVS}{IC Validator Explorer LVS}
\label{chap:explvs}

本章说明 IC Validator Explorer LVS 的功能。

IC Validator Explorer LVS 包括以下各节：
\begin{itemize}
  \item IC Validator Explorer LVS 简介
  \item 配置 Explorer LVS
  \item 执行 Explorer LVS
  \item Explorer LVS 许可证
\end{itemize}

% ============================================================
\section{IC Validator Explorer LVS 简介}
\subsection*{Introduction to IC Validator Explorer LVS}

在签核 LVS 早期找到 LVS 错误的根因，对于缩短整体签核周期、满足流片进度至关重要。Explorer LVS 面向片上系统（SOC）全芯片，提供可在较短周转时间与高效机器资源下定位 LVS 问题根因的方案，从而降低整体签核时间与工作量。
Explorer LVS 首先聚焦源于 SOC 集成的 LVS 问题，再逐步扩展分析，覆盖其他潜在 LVS 问题，例如来自器件/基础层或更低层次 cell 的 LVS 错误。
为检查 SOC 集成相关问题，Explorer LVS 默认仅对构成 SOC Top 的 equivalence cells 与金属层执行基于 text 的 short 检查（例如构建主 SOC 结构的 P\&R、PowerGrid、占位多边形 fill cells 等；相对地，宏与 IP 等则实例化在主 SOC 结构之下）。检查 LVS 问题的范围可按需通过 Explorer LVS 提供的灵活选项扩展。
Explorer LVS 有助于提升签核 LVS 的生产力与效率，但并不取代传统 LVS 工具。最终签核 LVS 仍须使用传统 LVS 工具运行。

% ============================================================
\section{配置 Explorer LVS}
\subsection*{Configuring Explorer LVS}

表~\ref{tab:explorer-lvs-cli} 描述可用于配置 Explorer LVS 的命令行选项。

\begin{longtable}{@{}>{\raggedright\arraybackslash\ttfamily}p{0.36\textwidth} p{0.56\textwidth}@{}}
\caption{Explorer LVS 命令行选项}\label{tab:explorer-lvs-cli}\\
\toprule
\textbf{\textrm{语法}} & \textbf{\textrm{说明}} \\
\midrule
\endfirsthead
\caption[]{Explorer LVS 命令行选项（续）}\\
\toprule
\textbf{\textrm{语法}} & \textbf{\textrm{说明}} \\
\midrule
\endhead
\bottomrule
\endfoot
-explorer lvs
  & 运行 LVS 的 Explorer 模式。对 \cmd{-explorer\_lvs\_cells}、\cmd{-explorer\_lvs\_layers}、\cmd{-explorer\_lvs\_checks}，分别采用 AUTO、METAL、SHORTS 默认值，并对 SOC Top cell 执行基于 text 的 short 检查（SOC Top cell 示例见「IC Validator Explorer LVS 简介」）。默认行为也可由本表后文所述这三个命令行选项更改。无论这些选项如何设置，Explorer LVS 均跳过器件提取。 \\
-explorer\_lvs\_cells [AUTO | ALL]
  & 可选。指定 Explorer LVS 要检查的 cells。\\
  \cmd{AUTO}（默认）：仅考虑与顶层相关的 cells。当 \cmd{-explorer\_lvs\_checks=SHORTS} 或 \cmd{PG\_SHORTS} 时，删除在 \cmd{equiv\_options()}、\cmd{lvs\_black\_box\_options()} 中定义的、或由 \cmd{lvs\_options(generate\_user\_equivs)} 生成的全部版图 cells。\\
  \cmd{ALL}：考虑全部 cells。 \\
-explorer\_lvs\_layers [METAL | ALL | user\_custom\_layer\_list]
  & 可选。指定 Explorer LVS 要检查的层。\\
  \cmd{METAL}（默认）：仅考虑金属层，或空的 FEOL 层（如器件端子层）。\\
  \cmd{ALL}：考虑全部层。\\
  \cmd{layer1}, \cmd{layer2}, \ldots{}：仅考虑金属层或空 FEOL 层，但排除 \cmd{layer1}、\cmd{layer2} 等所列层。 \\
-explorer\_lvs\_checks [SHORTS | ALL | PG\_SHORTS]
  & 可选。仅执行基于 text 的 short 检查，或同时执行基于 text 的 short 检查与 LVS compare。\\
  \cmd{SHORTS}（默认）：仅报告 \cmd{text\_short} 违例。\\
  \cmd{ALL}：报告 \cmd{text\_short} 违例与 LVS Compare 结果。\\
  \cmd{PG\_SHORTS}：报告电源与地的 \cmd{text\_short} 违例。\\
  PG shorts 定义：\cmd{<power, power>}、\cmd{<power, ground>}、\cmd{<ground, ground>}。\\
  non-PG shorts 定义：\cmd{<power, signal>}、\cmd{<ground, signal>}、\cmd{<signal, signal>}。 \\
-explorer\_lvs\_preserve\_cells [preserve\_cell\_list]
  & 可选。指定在 Explorer LVS 中应予保留、不被删除或 black-box 的 cells，以便将这些 cells 纳入 Explorer LVS。 \\
-explorer\_lvs\_preserve\_cell\_file "file path"
  & 指定 \cmd{-explorer\_lvs\_preserve\_cells} 的文件输入形式。每行一个待保留的 cell。 \\
-explorer\_lvs\_preserve\_hierarchy\_cells "space separated cellnames"|*
  & 指定在以 \cmd{-explorer lvs|lvs:1|lvs:2} 运行时，不应被删除或 black-box 的 \cmd{equiv\_cells} 及其子级 \cmd{equiv\_cells}。 \\
-explorer\_lvs\_preserve\_hierarchy\_cells\_file "file path"
  & 指定 \cmd{-explorer\_lvs\_preserve\_hierarchy\_cells} 的文件输入形式。每行一个 cell。 \\
\end{longtable}

Explorer LVS 命令行选项的全部可能用法组合如图~\ref{fig:explorer-lvs-cli-combos} 所示。

\figplaceholder{Figure 54: Combinations of Explorer LVS Command-Line Options}{Explorer LVS 命令行选项组合}{fig:explorer-lvs-cli-combos}

% ============================================================
\section{执行 Explorer LVS}
\subsection*{Executing Explorer LVS}

要执行 Explorer LVS，在命令行设置 \cmd{-explorer lvs}。该命令行选项对 \cmd{-explorer\_lvs\_cells}、\cmd{-explorer\_lvs\_layers}、\cmd{-explorer\_lvs\_checks} 分别采用 AUTO、METAL、SHORTS 默认值。Explorer LVS 仅对金属层上的 SOC Top cells 执行基于 text 的 short 检查。
默认行为可通过三个 Explorer LVS 选项 \cmd{-explorer\_lvs\_cells}、\cmd{-explorer\_lvs\_layers} 与 \cmd{-explorer\_lvs\_checks} 更改，见表~\ref{tab:explorer-lvs-cli}。
用 \cmd{-explorer\_lvs\_preserve\_cells} 指定的某些 cells 可在 Explorer LVS 中予以保留、不被删除或 black-box，从而纳入 Explorer LVS。
执行 Explorer LVS 的示例见表~\ref{tab:explorer-lvs-examples}。

\begin{longtable}{@{}p{0.22\textwidth} >{\raggedright\arraybackslash\ttfamily}p{0.38\textwidth} p{0.30\textwidth}@{}}
\caption{执行 Explorer LVS 的示例}\label{tab:explorer-lvs-examples}\\
\toprule
\textbf{\textrm{模式}} & \textbf{\textrm{命令行}} & \textbf{\textrm{说明}} \\
\midrule
\endfirsthead
\toprule
\textbf{\textrm{模式}} & \textbf{\textrm{命令行}} & \textbf{\textrm{说明}} \\
\midrule
\endhead
\bottomrule
\endfoot
默认
  & \% icv -explorer lvs \textbackslash\newline runsetfile
  & 仅对金属层上的 SOC Top cells 做基于 text 的 short 检查 \\
在默认基础上加入全部 cells
  & \% icv -explorer lvs \textbackslash\newline -explorer\_lvs\_cells ALL \textbackslash\newline runsetfile
  & 对全部 cells（仅金属层）做基于 text 的 short 检查 \\
在默认基础上加入全部层
  & \% icv -explorer lvs \textbackslash\newline -explorer\_lvs\_layers ALL \textbackslash\newline runsetfile
  & 对 SOC Top cells（全部层）做基于 text 的 short 检查 \\
在默认基础上加入全部 cells 与层
  & \% icv -explorer lvs \textbackslash\newline -explorer\_lvs\_cells ALL \textbackslash\newline -explorer\_lvs\_layers ALL \textbackslash\newline runsetfile
  & 对全部 cells、全部层做基于 text 的 short 检查 \\
在默认基础上加入 LVS compare 检查
  & \% icv -explorer lvs \textbackslash\newline -explorer\_lvs\_checks ALL \textbackslash\newline runsetfile
  & 对 SOC Top cells（仅金属层）做基于 text 的 short 检查，并附带 LVS compare \\
在 Explorer LVS 中保留 CellA 与 CellB
  & \% icv -explorer lvs \textbackslash\newline -explorer\_lvs\_preserve\_cells \textbackslash\newline "CellA CellB" runsetfile
  & 与默认行为相同，唯一区别是 Explorer LVS 保证保留 CellA 与 CellB \\
\end{longtable}

% ============================================================
\section{Explorer LVS 许可证}
\subsection*{Explorer LVS Licenses}

IC Validator Explorer LVS 需要 Elite 许可。该功能需要下列 IC Validator 许可证：
\begin{itemize}
  \item \cmd{ICValidator-Manager}
  \item \cmd{ICValidator2-GeometryEngine}
\end{itemize}

\begin{seeAlsoBox}
更多信息见本手册「许可证与资源需求」一章。
\end{seeAlsoBox}
